\chapter{Discussão}
\label{chap:discussao}

O capítulo de Discussão realiza a \textbf{interpretação crítica} dos resultados obtidos, conectando os achados empíricos às hipóteses iniciais e ao arcabouço teórico apresentado na Fundamentação Teórica e nos Trabalhos Relacionados.

Uma discussão científica de excelência deve abordar os seguintes pontos:
\begin{enumerate}
    \item \textbf{Resposta às Questões de Pesquisa:} Explicitar em que medida cada Questão de Pesquisa (QP1, QP2, QP3) foi respondida pelos dados.
    \item \textbf{Confrontação com a Literatura:} Comparar os achados com os estudos de referência (ex.: \enquote{Estes resultados corroboram as constatações de \cite{akhtar2018threat}, porém divergem quanto...}).
    \item \textbf{Explicação de Anomalias:} Analisar eventuais comportamentos atípicos ou cenários em que o desempenho ficou aquém do esperado, justificando as causas técnicas ou metodológicas subjacentes.
    \item \textbf{Implicações Práticas:} Destacar os desdobramentos operacionais e o valor gerado para a prática de engenharia de software e inovação.
\end{enumerate}